\chapter{Algebraic Bethe Ansatz}
\label{chap:algebraic-bethe-ansatz}

The coordinate Bethe ansatz solves the XXZ chain by constructing wave functions in
ordered magnon-coordinate sectors.  The algebraic Bethe ansatz solves the same
problem from the transfer-matrix side.  Instead of writing a many-body wave
function explicitly, one uses the Yang--Baxter equation to build an operator
algebra.  Bethe states are then obtained by applying a creation operator
$B(u)$, extracted from the monodromy matrix, to a simple reference state.

The conceptual chain is
\[
\begin{gathered}
    R\text{-matrix satisfying YBE}
    \Longrightarrow
    \text{RTT algebra for }A(u),B(u),C(u),D(u) \\
    \Longrightarrow
    \text{Bethe vectors }B(\lambda_1)\cdots B(\lambda_M)\ket{0} \\
    \Longrightarrow
    \text{transfer-matrix eigenvalues and Bethe equations}.
\end{gathered}
\]
Thus the algebraic Bethe ansatz diagonalizes not only one Hamiltonian, but the
whole commuting transfer-matrix family.  The Hamiltonian eigenvalues are obtained
afterwards by taking the logarithmic derivative of the transfer-matrix
eigenvalue.

Throughout this chapter we use the homogeneous six-vertex convention of
\cref{chap:six-vertex-xxz}.  The local $R$-matrix is
\begin{equation}
\label{eq:aba-six-vertex-R}
    R(u)=
    \begin{pmatrix}
        \sinh(u+\eta)&0&0&0\\
        0&\sinh u&\sinh\eta&0\\
        0&\sinh\eta&\sinh u&0\\
        0&0&0&\sinh(u+\eta)
    \end{pmatrix},
    \qquad
    \Delta=\cosh\eta .
\end{equation}
The trigonometric critical regime follows by analytic continuation
$\eta\mapsto \ii\gamma$, giving $\Delta=\cos\gamma$.  We keep the hyperbolic
notation because it makes the algebraic formulas compact.

\section{Monodromy matrix}
\label{sec:aba-monodromy-matrix}

Let $V_a\simeq\CC^2$ be an auxiliary spin-$1/2$ space and
\begin{equation}
\label{eq:aba-quantum-space}
    \calH_L=V_1\otimes\cdots\otimes V_L,
    \qquad V_j\simeq\CC^2
\end{equation}
be the quantum spin-chain Hilbert space.  The monodromy matrix is the ordered
product
\begin{equation}
\label{eq:aba-monodromy-definition}
    M_a(u)
    =R_{aL}(u)R_{a,L-1}(u)\cdots R_{a1}(u).
\end{equation}
It is a $2\times2$ matrix in the auxiliary space, whose entries are operators on
$\calH_L$:
\begin{equation}
\label{eq:aba-monodromy-ABCD}
    M_a(u)=
    \begin{pmatrix}
        A(u)&B(u)\\
        C(u)&D(u)
    \end{pmatrix}_a .
\end{equation}
The transfer matrix is the auxiliary trace
\begin{equation}
\label{eq:aba-transfer-matrix}
    \tau(u)=\Tr_a M_a(u)=A(u)+D(u).
\end{equation}
The entries $A,B,C,D$ are not four independent Hamiltonians.  They are
noncommuting operator-valued functions constrained by the RTT algebra below.

It is useful to write the local $R$-operator as a matrix in the auxiliary space
with quantum-site operators as entries.  In the convention of
\cref{eq:aba-six-vertex-R},
\begin{equation}
\label{eq:aba-local-L-operator}
    R_{aj}(u)=
    \begin{pmatrix}
        a_R(u)e^{11}_j+b_R(u)e^{22}_j
        &c_R e^{21}_j\\
        c_R e^{12}_j
        &b_R(u)e^{11}_j+a_R(u)e^{22}_j
    \end{pmatrix}_a,
\end{equation}
where
\begin{equation}
\label{eq:aba-local-weights}
    a_R(u)=\sinh(u+\eta),
    \qquad
    b_R(u)=\sinh u,
    \qquad
    c_R=\sinh\eta,
\end{equation}
and $e^{\alpha\beta}_j=\ket{\alpha}_j\bra{\beta}_j$.  With
$\ket{\uparrow}=\ket{1}$ and $\ket{\downarrow}=\ket{2}$, one has
$e^{21}=\sigma^-$ and $e^{12}=\sigma^+$.

Choose the ferromagnetic reference state
\begin{equation}
\label{eq:aba-reference-state}
    \ket{0}=\ket{\uparrow\uparrow\cdots\uparrow} .
\end{equation}
On this state every local $R_{aj}(u)$ becomes upper triangular in auxiliary
space:
\begin{equation}
\label{eq:aba-local-triangular-action}
    R_{aj}(u)\ket{\uparrow}_j
    =
    \begin{pmatrix}
        a_R(u)\ket{\uparrow}_j & c_R\ket{\downarrow}_j\\
        0&b_R(u)\ket{\uparrow}_j
    \end{pmatrix}_a .
\end{equation}
Multiplying over all sites preserves triangularity.  Hence
\begin{equation}
\label{eq:aba-pseudovacuum-action}
    C(u)\ket{0}=0,
    \qquad
    A(u)\ket{0}=\alpha(u)\ket{0},
    \qquad
    D(u)\ket{0}=\delta(u)\ket{0},
\end{equation}
with vacuum eigenvalues
\begin{equation}
\label{eq:aba-vacuum-eigenvalues}
    \alpha(u)=a_R(u)^L=\sinh^L(u+\eta),
    \qquad
    \delta(u)=b_R(u)^L=\sinh^L u .
\end{equation}
The operator $B(u)$ lowers one or more spins when expanded as an operator on the
chain.  Its leading action on the reference state is a one-magnon superposition,
so $B(u)$ is the algebraic creation operator for Bethe quasiparticles.

\begin{remark}[Pseudovacuum]
The word ``vacuum'' here is algebraic.  For the antiferromagnetic XXZ chain the
physical ground state is usually not $\ket{0}$.  The state $\ket{0}$ is used
because it is a highest-weight state for the RTT algebra: $C(u)$ annihilates it,
while $A(u)$ and $D(u)$ act diagonally.
\end{remark}

\section{RTT relation and commuting transfer matrices}
\label{sec:aba-rtt-commuting-transfer}

The Yang--Baxter equation for the six-vertex $R$-matrix implies the local
intertwining relation
\begin{equation}
\label{eq:aba-local-RLL}
    R_{ab}(u-v)R_{aj}(u)R_{bj}(v)
    =
    R_{bj}(v)R_{aj}(u)R_{ab}(u-v).
\end{equation}
Multiplying this identity over all sites gives the RTT relation
\begin{equation}
\label{eq:aba-RTT-relation}
    R_{ab}(u-v)M_a(u)M_b(v)
    =
    M_b(v)M_a(u)R_{ab}(u-v).
\end{equation}
This is the defining algebraic relation of the algebraic Bethe ansatz.  It is an
operator identity on
\begin{equation}
    V_a\otimes V_b\otimes\calH_L .
\end{equation}
The two auxiliary spaces $a$ and $b$ are bookkeeping devices; the nontrivial
operators act on the same quantum chain $\calH_L$.

\begin{theorem}[Commuting transfer matrices]
\label{thm:aba-commuting-transfer-matrices}
The RTT relation implies
\begin{equation}
\label{eq:aba-transfer-commutativity}
    [\tau(u),\tau(v)]=0,
    \qquad
    \text{for all }u,v .
\end{equation}
\end{theorem}

\begin{proof}
Take $\Tr_a\Tr_b$ of \cref{eq:aba-RTT-relation}.  By cyclicity of the trace in
the auxiliary spaces, $R_{ab}(u-v)$ can be moved through the trace and canceled.
This gives
\begin{equation}
    \Tr_a M_a(u)\,\Tr_b M_b(v)
    =
    \Tr_b M_b(v)\,\Tr_a M_a(u),
\end{equation}
which is exactly \cref{eq:aba-transfer-commutativity}.
\end{proof}

The commutativity is the spectral reason for integrability.  Expanding around the
regular point gives mutually commuting charges,
\begin{equation}
\label{eq:aba-conserved-charges-from-transfer}
    \log\tau(u)=\sum_{n\geq0}u^n Q_n,
    \qquad
    [Q_m,Q_n]=0 .
\end{equation}
The XXZ Hamiltonian is one of these charges:
\begin{equation}
\label{eq:aba-xxz-from-log-transfer}
    H_{\rm XXZ}^{(+)}
    =
    2J\sinh\eta
    \left.\frac{\dd}{\dd u}\log\tau(u)\right|_{u=0}
    -2J\cosh\eta\,L,
\end{equation}
which is the same normalization as
\begin{equation}
\label{eq:aba-xxz-hamiltonian}
    H_{\rm XXZ}^{(+)}
    =
    J\sum_{j=1}^{L}
    \left(
        \sigma_j^x\sigma_{j+1}^x
        +\sigma_j^y\sigma_{j+1}^y
        +\Delta(\sigma_j^z\sigma_{j+1}^z-1)
    \right),
    \qquad
    \Delta=\cosh\eta .
\end{equation}
The constant fixes the energy of the fully polarized reference state to zero.

\section{The fundamental commutation relations}
\label{sec:aba-fundamental-commutation-relations}

The RTT relation contains many component equations.  The algebraic Bethe ansatz
only needs the relations that move $A(u)$ and $D(u)$ through products of $B$'s.
Introduce
\begin{equation}
\label{eq:aba-f-and-g-functions}
    f(u)=\frac{\sinh(u+\eta)}{\sinh u},
    \qquad
    g(u)=\frac{\sinh\eta}{\sinh u} .
\end{equation}
With the monodromy order \cref{eq:aba-monodromy-definition}, the needed RTT
relations are
\begin{align}
\label{eq:aba-BB-commutation}
    B(u)B(v)&=B(v)B(u),\\
\label{eq:aba-AB-commutation}
    A(u)B(v)
    &=
    f(v-u)B(v)A(u)-g(v-u)B(u)A(v),\\
\label{eq:aba-DB-commutation}
    D(u)B(v)
    &=
    f(u-v)B(v)D(u)-g(u-v)B(u)D(v).
\end{align}
The first identity says that Bethe creation operators commute.  The second and
third identities have two parts.  The first terms preserve the set of Bethe
rapidities and will produce the transfer-matrix eigenvalue.  The second terms
replace one rapidity $v$ by the spectral parameter $u$; these are the unwanted
terms whose cancellation gives the Bethe equations.

\begin{remark}[Convention dependence]
If one reverses the monodromy ordering, changes $R(u-v)$ to $R(v-u)$, or shifts
the spectral parameter by $\eta/2$, the displayed formulas may be inverted or
shifted.  The invariant content is the same: the ratio of vacuum eigenvalues is
balanced against the product of two-body scattering factors.
\end{remark}

\section{Bethe vectors}
\label{sec:aba-bethe-vectors}

For rapidities $\lambda_1,
\ldots,\lambda_M$, define the Bethe vector
\begin{equation}
\label{eq:aba-bethe-vector}
    \ket{\lambda_1,\ldots,\lambda_M}
    =
    B(\lambda_1)B(\lambda_2)\cdots B(\lambda_M)\ket{0} .
\end{equation}
Because the $B$'s commute, this vector is symmetric under permutations of the
rapidities.  It lies in the sector with $M$ down spins:
\begin{equation}
\label{eq:aba-magnetization-sector}
    S^z_{\rm tot}\ket{\lambda_1,\ldots,\lambda_M}
    =
    \left(\frac{L}{2}-M\right)
    \ket{\lambda_1,\ldots,\lambda_M} .
\end{equation}
This matches the coordinate Bethe ansatz, where $M$ was the number of magnons.

\begin{example}[One Bethe excitation]
For $M=1$,
\begin{equation}
    \ket{\lambda}=B(\lambda)\ket{0} .
\end{equation}
Using \cref{eq:aba-AB-commutation,eq:aba-DB-commutation},
\begin{equation}
\label{eq:aba-one-particle-action}
\begin{aligned}
    \tau(u)B(\lambda)\ket{0}
    ={}&
    \left[
        \alpha(u)f(\lambda-u)+\delta(u)f(u-\lambda)
    \right]B(\lambda)\ket{0}
    \\
    &
    +g(\lambda-u)
    \left[
        \delta(\lambda)-\alpha(\lambda)
    \right]B(u)\ket{0} .
\end{aligned}
\end{equation}
The first line is proportional to the original vector.  The second line is
unwanted.  It vanishes exactly when
\begin{equation}
\label{eq:aba-one-particle-bethe-equation}
    \alpha(\lambda)=\delta(\lambda),
    \qquad
    \left[\frac{\sinh(\lambda+\eta)}{\sinh\lambda}\right]^L=1 .
\end{equation}
This is the one-magnon momentum quantization condition in algebraic form.
\end{example}

The many-particle calculation is the repeated version of the same mechanism.
Move $A(u)$ and $D(u)$ through the string of $B$'s until they reach the
pseudovacuum.  The wanted part keeps all rapidities fixed.  Each unwanted term
replaces one $\lambda_j$ by $u$.

\section{Transfer-matrix eigenvalue and Bethe equations}
\label{sec:aba-transfer-eigenvalue-bethe-equations}

Acting with $\tau(u)=A(u)+D(u)$ on the Bethe vector gives the off-shell form
\begin{equation}
\label{eq:aba-off-shell-action}
    \tau(u)\ket{\lambda_1,
\ldots,\lambda_M}
    =
    \Lambda(u\mid\boldsymbol{\lambda})
    \ket{\lambda_1,
\ldots,\lambda_M}
    +
    \sum_{j=1}^{M}
    \mathcal{F}_j(u\mid\boldsymbol{\lambda})
    B(u)\prod_{\ell\neq j}B(\lambda_\ell)\ket{0} .
\end{equation}
Here the wanted coefficient is
\begin{equation}
\label{eq:aba-transfer-eigenvalue}
\boxed{
\begin{aligned}
    \Lambda(u\mid\boldsymbol{\lambda})
    ={}&
    \alpha(u)
    \prod_{j=1}^{M}f(\lambda_j-u)
    +
    \delta(u)
    \prod_{j=1}^{M}f(u-\lambda_j)
    \\
    ={}&
    \sinh^L(u+\eta)
    \prod_{j=1}^{M}
    \frac{\sinh(\lambda_j-u+\eta)}{\sinh(\lambda_j-u)}
    \\
    &+
    \sinh^L u
    \prod_{j=1}^{M}
    \frac{\sinh(u-\lambda_j+\eta)}{\sinh(u-\lambda_j)} .
\end{aligned}}
\end{equation}
Up to a nonzero scalar factor, the unwanted coefficient is
\begin{equation}
\label{eq:aba-unwanted-coefficient}
    \mathcal{F}_j(u\mid\boldsymbol{\lambda})
    \propto
    \delta(\lambda_j)
    \prod_{\ell\neq j}f(\lambda_j-\lambda_\ell)
    -
    \alpha(\lambda_j)
    \prod_{\ell\neq j}f(\lambda_\ell-\lambda_j) .
\end{equation}
Therefore the Bethe vector is an eigenvector of $\tau(u)$ for all $u$ if
\begin{equation}
\label{eq:aba-bethe-equations-abstract}
    \frac{\alpha(\lambda_j)}{\delta(\lambda_j)}
    =
    \prod_{\substack{\ell=1\\ \ell\neq j}}^{M}
    \frac{f(\lambda_j-\lambda_\ell)}{f(\lambda_\ell-\lambda_j)},
    \qquad
    j=1,
\ldots,M .
\end{equation}
Substituting the six-vertex functions gives the algebraic Bethe equations
\begin{equation}
\label{eq:aba-xxz-bethe-equations}
\boxed{
    \left[
    \frac{\sinh(\lambda_j+\eta)}{\sinh\lambda_j}
    \right]^L
    =
    \prod_{\substack{\ell=1\\ \ell\neq j}}^{M}
    \frac{\sinh(\lambda_j-\lambda_\ell+\eta)}
         {\sinh(\lambda_j-\lambda_\ell-\eta)} ,
    \qquad
    j=1,
\ldots,M .
}
\end{equation}
When these equations hold, all unwanted terms in \cref{eq:aba-off-shell-action}
vanish, and
\begin{equation}
\label{eq:aba-on-shell-transfer-action}
    \tau(u)\ket{\lambda_1,
\ldots,\lambda_M}
    =
    \Lambda(u\mid\boldsymbol{\lambda})
    \ket{\lambda_1,
\ldots,\lambda_M} .
\end{equation}

\begin{principle}[Bethe equations as analyticity conditions]
The transfer-matrix eigenvalue \cref{eq:aba-transfer-eigenvalue} appears to have
simple poles at $u=\lambda_j$.  The Bethe equations are exactly the conditions
that the residues cancel.  Thus the same equations can be read either as
cancellation of unwanted vectors or as analyticity of the transfer-matrix
eigenvalue.
\end{principle}

\section{Hamiltonian eigenvalues}
\label{sec:aba-hamiltonian-eigenvalues}

The transfer-matrix eigenvalue contains the eigenvalues of all commuting charges.
Using \cref{eq:aba-xxz-from-log-transfer}, the XXZ energy is
\begin{equation}
\label{eq:aba-energy-from-transfer-eigenvalue}
    E(\boldsymbol{\lambda})
    =
    2J\sinh\eta
    \left.
    \frac{\dd}{\dd u}
    \log\Lambda(u\mid\boldsymbol{\lambda})
    \right|_{u=0}
    -2J\cosh\eta\,L .
\end{equation}
For the eigenvalue \cref{eq:aba-transfer-eigenvalue}, this gives
\begin{equation}
\label{eq:aba-energy-rapidity-form}
\boxed{
    E(\boldsymbol{\lambda})
    =
    2J\sinh^2\eta
    \sum_{j=1}^{M}
    \frac{1}{\sinh\lambda_j\,\sinh(\lambda_j+\eta)} .
}
\end{equation}
The empty Bethe state $M=0$ has $E=0$, as required by the constant chosen in
\cref{eq:aba-xxz-hamiltonian}.

The lattice momentum is obtained from the regular-point transfer matrix.  Since
$R(0)=\sinh\eta\,P$, one has
\begin{equation}
\label{eq:aba-transfer-at-zero-translation}
    \tau(0)=(\sinh\eta)^L U,
\end{equation}
where $U$ is the one-site translation operator, with the orientation determined
by the monodromy ordering.  Therefore
\begin{equation}
\label{eq:aba-momentum-from-roots}
    U\ket{\boldsymbol{\lambda}}
    =
    \left[
    \prod_{j=1}^{M}
    \frac{\sinh(\lambda_j+\eta)}{\sinh\lambda_j}
    \right]
    \ket{\boldsymbol{\lambda}} .
\end{equation}
Depending on whether one defines $U$ as a left or right shift, this eigenvalue is
$\ee^{\ii P}$ or $\ee^{-\ii P}$.  This is only an orientation convention.

If one introduces a one-particle momentum by
\begin{equation}
\label{eq:aba-momentum-rapidity-map}
    \ee^{\ii k(\lambda)}
    =
    \frac{\sinh(\lambda+\eta)}{\sinh\lambda},
\end{equation}
then
\begin{equation}
\label{eq:aba-energy-momentum-form}
    \frac{2J\sinh^2\eta}{\sinh\lambda\,\sinh(\lambda+\eta)}
    =
    4J\bigl(\cos k(\lambda)-\cosh\eta\bigr),
\end{equation}
so \cref{eq:aba-energy-rapidity-form} reproduces the coordinate Bethe-ansatz
energy
\begin{equation}
\label{eq:aba-coordinate-energy-recovered}
    E=4J\sum_{j=1}^{M}\bigl(\cos k_j-\Delta\bigr).
\end{equation}
If the opposite translation orientation is used, $k$ is replaced by $-k$ and the
energy is unchanged.

\section[Symmetric rapidities and CBA]{Symmetric rapidities and comparison with the coordinate Bethe ansatz}
\label{sec:aba-comparison-coordinate-bethe-ansatz}

The rapidity in \cref{eq:aba-xxz-bethe-equations} is tied to the unshifted
six-vertex $R$-matrix.  A more symmetric variable is
\begin{equation}
\label{eq:aba-symmetric-rapidity}
    \mu=\lambda+\frac{\eta}{2} .
\end{equation}
Then
\begin{equation}
\label{eq:aba-symmetric-momentum-map}
    \ee^{\ii k(\mu)}
    =
    \frac{\sinh(\mu+\eta/2)}{\sinh(\mu-\eta/2)},
\end{equation}
and the Bethe equations become
\begin{equation}
\label{eq:aba-symmetric-bethe-equations}
    \left[
    \frac{\sinh(\mu_j+\eta/2)}{\sinh(\mu_j-\eta/2)}
    \right]^L
    =
    \prod_{\substack{\ell=1\\ \ell\neq j}}^{M}
    \frac{\sinh(\mu_j-\mu_\ell+\eta)}
         {\sinh(\mu_j-\mu_\ell-\eta)} .
\end{equation}
After analytic continuation $\eta=\ii\gamma$, this is the standard critical XXZ
form
\begin{equation}
\label{eq:aba-critical-xxz-bethe-equations}
    \left[
    \frac{\sinh(\mu_j+\ii\gamma/2)}{\sinh(\mu_j-\ii\gamma/2)}
    \right]^L
    =
    \prod_{\substack{\ell=1\\ \ell\neq j}}^{M}
    \frac{\sinh(\mu_j-\mu_\ell+\ii\gamma)}
         {\sinh(\mu_j-\mu_\ell-\ii\gamma)} .
\end{equation}
This is the same equation obtained in the coordinate Bethe ansatz, up to the
orientation convention for lattice momentum.

The two approaches differ in what they make manifest:
\begin{center}
\begin{tabular}{@{}L{0.30\linewidth}L{0.30\linewidth}L{0.30\linewidth}@{}}
\toprule
Method & Main calculation & Output \\
\midrule
Coordinate Bethe ansatz & Solve contact conditions in real-space wave functions & Two-body scattering phase and momentum quantization \\
Algebraic Bethe ansatz & Commute $A(u)$ and $D(u)$ through products of $B$'s & Transfer-matrix eigenvalue and Bethe equations \\
\bottomrule
\end{tabular}
\end{center}
The coordinate method emphasizes factorized scattering.  The algebraic method
explains why factorized scattering exists: the RTT relation is the operator form
of the Yang--Baxter equation.

\section{Higher conserved charges}
\label{sec:aba-higher-conserved-charges}

Because the Bethe vector is an eigenvector of $\tau(u)$ for all $u$, it is also a
simultaneous eigenvector of all charges generated by $\log\tau(u)$.  If
\begin{equation}
\label{eq:aba-log-transfer-charge-expansion}
    \log\tau(u)=\sum_{n\geq0}u^n Q_n,
\end{equation}
then
\begin{equation}
\label{eq:aba-charge-eigenvalues}
    Q_n\ket{\boldsymbol{\lambda}}
    =q_n(\boldsymbol{\lambda})\ket{\boldsymbol{\lambda}},
    \qquad
    q_n(\boldsymbol{\lambda})
    =
    \left.
    \frac{1}{n!}
    \frac{\dd^n}{\dd u^n}
    \log\Lambda(u\mid\boldsymbol{\lambda})
    \right|_{u=0}
\end{equation}
up to the same normalization convention used in defining the charges.  The
Hamiltonian is the first nontrivial local charge.  Higher $Q_n$ have larger
spatial support and provide the extensive set of conservation laws characteristic
of an integrable chain.

\begin{remark}[Why this is stronger than diagonalizing one Hamiltonian]
A generic many-body Hamiltonian may be diagonalized in principle, but its
eigenvectors do not usually diagonalize a whole analytic family of local and
quasilocal conserved charges.  The algebraic Bethe ansatz gives simultaneous
eigenvectors of the transfer matrix for all spectral parameters.  This is the
finite-size algebraic signature of integrability.
\end{remark}

\section{Trigonometric and rational limits}
\label{sec:aba-trigonometric-rational-limits}

For $\abs{\Delta}<1$, set
\begin{equation}
\label{eq:aba-eta-to-gamma}
    \eta=\ii\gamma,
    \qquad
    0<\gamma<\pi,
    \qquad
    \Delta=\cos\gamma .
\end{equation}
The same algebraic derivation applies, with $\sinh$ functions analytically
continued.  In the symmetric rapidity convention, the Bethe equations are exactly
\cref{eq:aba-critical-xxz-bethe-equations}.

The isotropic XXX chain is the rational degeneration.  Scaling the spectral
parameter and anisotropy together produces the rational $R$-matrix
\begin{equation}
\label{eq:aba-rational-R-schematic}
    R_{\rm XXX}(u)=u\,\id+\ii P,
    \qquad
    \text{up to an overall scalar normalization.}
\end{equation}
In the usual XXX rapidity variable, the Bethe equations become
\begin{equation}
\label{eq:aba-xxx-bethe-equations}
    \left(
    \frac{\lambda_j+\ii/2}{\lambda_j-\ii/2}
    \right)^L
    =
    \prod_{\substack{\ell=1\\ \ell\neq j}}^{M}
    \frac{\lambda_j-\lambda_\ell+\ii}
         {\lambda_j-\lambda_\ell-\ii} .
\end{equation}
This rational limit will reappear when the XXX chain is discussed as a special
case of the XXZ chain.

\section{Common pitfalls}
\label{sec:aba-common-pitfalls}

\begin{remark}[The operators $A,B,C,D$ are nonlocal]
Although $R_{aj}(u)$ is local in the quantum site $j$, the monodromy entries
$A(u),B(u),C(u),D(u)$ are products over the whole chain.  The operator $B(u)$ is
therefore not a local spin-lowering operator.  It is a spectral-parameter
dependent, nonlocal creation operator for Bethe quasiparticles.
\end{remark}

\begin{remark}[Bethe roots cannot sit at generic poles]
The formulas above contain denominators such as $\sinh\lambda_j$ and
$\sinh(\lambda_j-\lambda_\ell)$.  Ordinary Bethe roots are assumed to avoid these
poles.  Singular solutions can occur at special points and require a limiting
procedure.  They are not part of the basic algebraic derivation.
\end{remark}

\begin{remark}[Completeness is a separate question]
The algebraic Bethe ansatz constructs eigenvectors whenever the Bethe equations
hold and the resulting vector is nonzero.  Showing that all eigenvectors are
obtained requires a separate completeness analysis.  For the finite XXZ chain
this is well understood, but it is conceptually distinct from deriving the Bethe
equations.
\end{remark}

\section{Summary}
\label{sec:aba-summary}

The algebraic Bethe ansatz for the XXZ chain can be compressed into the following
sequence:
\begin{equation}
\label{eq:aba-summary-chain}
\boxed{
\begin{gathered}
    R(u)\text{ satisfies YBE}
    \Longrightarrow
    R_{ab}(u-v)M_a(u)M_b(v)=M_b(v)M_a(u)R_{ab}(u-v),
    \\
    M_a(u)=
    \begin{pmatrix}A(u)&B(u)\\ C(u)&D(u)\end{pmatrix}_a,
    \qquad
    \tau(u)=A(u)+D(u),
    \\
    \ket{\boldsymbol{\lambda}}
    =B(\lambda_1)\cdots B(\lambda_M)\ket{0},
    \\
    \tau(u)\ket{\boldsymbol{\lambda}}
    =\Lambda(u\mid\boldsymbol{\lambda})\ket{\boldsymbol{\lambda}}
    \quad
    \text{if }\boldsymbol{\lambda}\text{ satisfies the Bethe equations},
    \\
    \left[
    \frac{\sinh(\lambda_j+\eta)}{\sinh\lambda_j}
    \right]^L
    =
    \prod_{\ell\neq j}
    \frac{\sinh(\lambda_j-\lambda_\ell+\eta)}
         {\sinh(\lambda_j-\lambda_\ell-\eta)} .
\end{gathered}}
\end{equation}
The Bethe equations arise from cancellation of unwanted terms in the RTT algebra.
The transfer-matrix eigenvalue then gives the complete set of commuting-charge
eigenvalues, including the XXZ Hamiltonian energy.  This is the algebraic
counterpart of the factorized scattering picture developed in the coordinate
Bethe-ansatz chapter.
